\section{Combinatorial preliminaries}
\label{chapter_combinatorics}

We recall only the notation used below. A partition of $n$ is a finite,
weakly decreasing sequence $\lambda=(\lambda_1,\lambda_2,\ldots,\lambda_k)$
of positive integers with $\sum_i\lambda_i=n$. We write $\lambda\leq
\mu$ in dominance order if
\[
  \sum_{i=1}^k\lambda_i\leq\sum_{i=1}^k\mu_i
  \qquad\text{for every }k\geq 1.
\]
The transpose partition, obtained by reflecting the Young diagram about
the main diagonal, is denoted by $\lambda^T$. A standard Young tableau
of shape $\lambda$ is a filling of its Young diagram by $1,\ldots,n$
whose entries increase from left to right and from bottom to top. We
denote the set of such tableaux by $\mathcal{Y}_\lambda$.

The Robinson--Schensted correspondence
\cite{robinson1938representations,schensted1961longest} is a bijection
\[
  \mathrm{RS}:S_n\xrightarrow{\sim}
  \bigsqcup_{\lambda\vdash n}
  \mathcal{Y}_\lambda\times\mathcal{Y}_\lambda,
  \qquad w\longmapsto(P_w,Q_w).
\]
We use the convention in which tableaux increase from bottom to top.
The only general property needed later is the following symmetry.

\begin{proposition}[pps:]{\cite{schutzenberger1963quelques}}
  If $\mathrm{RS}(w)=(P_w,Q_w)$, then
  $\mathrm{RS}(w^{-1})=(Q_w,P_w)$.
\end{proposition}

For partitions $\lambda,\mu\vdash n$, let
$K_{\lambda,\mu}(q)$ denote the Kostka--Foulkes polynomial. For
$T\in\mathcal{Y}_\lambda$, let $\operatorname{read}(T)$ be obtained by
reading each row from left to right, starting with the top row, and set
\[
  \operatorname{charge}(T)
  =\sum_{\substack{1\leq i<n\\
      i+1\text{ occurs to the right of }i\text{ in }\operatorname{read}(T)}}
    (n-i).
\]
With this convention \cite[Chapter III, \S6]{macdonald1995symmetric},
\[
  K_{\lambda,(1^n)}(q)
  =\sum_{T\in\mathcal{Y}_\lambda}q^{\operatorname{charge}(T)}.
\]
In particular,
\[
  K_{\lambda,(1^n)}(1)=|\mathcal{Y}_\lambda|.
\]
This convention agrees with the normalization used in
\Cref{thm:IH_kostka}.
